\section{晚唐的逐年压力：不把断裂压成一个年份}
\label{sec:annual-late-tang-rhythms}

779年至790年，德宗即位、两税法、河北军镇拒命、泾原兵变、奉天复都、平凉会盟和安西北庭联络终结相继发生。把这十余年压成“两税法以后藩镇坐大”，会看不见每一年里朝廷都在重新决定赦宥、征收、出兵和外交的先后。年度编年锚点的作用正在于保存这种尚未被大事件吸收的时间：它提示研究者回到当年本纪、诏令、会要和地方材料，而不把一条年序入口误写成已经解释完毕的社会事实。%
\footnote{这一段的逐年检索入口包括新增账本 V03-779-ANNAL-CENTRAL、V03-780-INSTITUTION-SOCIETY、V03-781-CROSS-POLITY、V03-783-ANNAL-CENTRAL、V03-784-INSTITUTION-SOCIETY、V03-785-ANNAL-CENTRAL、V03-787-CROSS-POLITY、V03-790-ANNAL-CENTRAL；具体节点见 LT-779-019 至 LT-790-027。年度锚点只锁定核查路径，不能代替原始条目。}

801年至829年也不应只留下“元和中兴”一个标签。西川的边防、永贞改革的短暂推进、刘辟与吴元济的叛乱、蔡州夜袭、佛骨争议、科举案、唐蕃会盟和沧景军镇的继承危机，使军政、财政、宗教和取士在不同年份以不同速度发生摩擦。宪宗朝能够在某些地区恢复军事主动，并不使盐税、转运、宦官和地方军人从制度中退出；穆宗以后军镇的再起，亦不表示每一项中央文书都已无效。%
\footnote{年度入口可参 V03-801-ANNAL-CENTRAL、V03-805-INSTITUTION-SOCIETY、V03-809-ANNAL-CENTRAL、V03-813-CROSS-POLITY、V03-815-ANNAL-CENTRAL、V03-817-ANNAL-CENTRAL、V03-819-INSTITUTION-SOCIETY、V03-821-ANNAL-CENTRAL、V03-824-ANNAL-CENTRAL、V03-828-ANNAL-CENTRAL、V03-829-INSTITUTION-SOCIETY；具体事件见 LT-801-028 至 LT-829-043。正史和《通鉴》的年序需同碑志、文书和制度志分层对读。}

835年至859年，甘露之变、回鹘汗国覆亡后的迁徙、吐蕃内部危机、会昌废佛、归义军的形成与宣宗末年的继承，构成一系列并不同步的区域转折。宫廷里禁军和宦官的关系、河西的地方军府、北边部众的安置、寺院财产的处置，不能由同一个“晚唐衰败”公式说明。逐年记录能防止我们把846年的新皇帝或851年的册命误读成全帝国一齐回稳的证明。%
\footnote{年度入口包括 V03-835-ANNAL-CENTRAL、V03-840-CROSS-POLITY、V03-842-CROSS-POLITY、V03-845-INSTITUTION-SOCIETY、V03-848-ANNAL-CENTRAL、V03-851-INSTITUTION-SOCIETY、V03-859-ANNAL-CENTRAL；具体事件见 LT-835-044 至 LT-859-053。回鹘、吐蕃与归义军的材料语言和保存地点不同，不能以中原本纪单线叙述。}

863年至907年，安南反复失守和收复、庞勋起事、王仙芝与黄巢的移动、成都行在、盐池冲突、军府在成都凤翔淮南的成形，以及洛阳朝廷最后的清洗，构成更密集的年度危机。它不是一场从876年直接通向907年的直线崩塌：港市、山地、军镇、禁军和行在各有短暂恢复和新的断裂。保留年度锚点，是为了让后续叙事继续追问每一年哪些税源、道路和军队仍可被调动，而不是让唐亡的结果提前吞没所有当时的选择。%
\footnote{年度入口可参 V03-863-ANNAL-CENTRAL、V03-866-ANNAL-CENTRAL、V03-868-INSTITUTION-SOCIETY、V03-873-ANNAL-CENTRAL、V03-874-ANNAL-CENTRAL、V03-879-CROSS-POLITY、V03-880-ANNAL-CENTRAL、V03-885-INSTITUTION-SOCIETY、V03-891-ANNAL-CENTRAL、V03-900-ANNAL-CENTRAL、V03-903-ANNAL-CENTRAL、V03-904-ANNAL-CENTRAL、V03-905-ANNAL-CENTRAL、V03-907-ANNAL-CENTRAL；具体事件见 LT-863-054 至 LT-907-080。}

\subsection{八七四至九〇七：流动作战中的逐年选择}

八七四年王仙芝在濮、曹聚众，连年灾害、赋役与盐禁冲突使流民、私盐贩和失业者能够跨州移动；八七九年黄巢军进入广州，港市与仓储遭受冲击后北返；八八〇年攻入长安称帝。年度顺序说明，起事者并非从山东直接抵达京师，而是在中原、岭南与关中不同的粮源、道路和防线间改变方向。广州的驻留、杀掠与贸易破坏的数字主要出自敌对叙事，长安居民的反应也不能由复国史书一概而论；但黄巢在关中面对粮食、官僚与防务的难题，正是大齐政权难以把军事流动转为稳定统治的条件。

八八一年僖宗至成都，西川的粮仓、城防和军政资源维持行在的诏令与调兵；八八二年朱温降唐受授宣武，八八三年唐、沙陀等军收复长安，军功与授官随即使军镇积累独立的兵粮和政治资本。八八五年王重荣与田令孜围绕调兵、盐利冲突，皇帝再出奔；八九一年王建取得成都，避难基地转为可独立经营的区域中心。由此看，收复长安不是恢复一切的年份，而是关中、西川、河东和汴州各自掌握资源的关系被重新排列的年份。

九〇三年朱温诛杀韩全诲等宦官并重编神策军，九〇四年迫昭宗迁洛阳且弑之，九〇五年在白马驿清除朝臣，九〇七年遂迫哀帝禅位。每一节点都缩小了唐室能依赖的禁军、旧都和官僚空间；也不能将仪式文本误作没有军事胁迫的“和平转移”。《旧唐书》本纪、朱全忠传与《资治通鉴》可校事件年序，神策军是否完全解体、白马驿死者与洛阳官员去向仍须用不同层次的材料分别核验。
